\section{Висновки}\label{sec:conclusions}

В результаті виконання курсової роботи розроблено та верифіковано в симуляції систему автономної інспекції складу роєм БПЛА з потоковою IoT-телеметрією та алгоритмом естафетної передачі місії.

\subsection*{Досягнуті результати}

\begin{enumerate}[itemsep=4pt]
  \item \textbf{Архітектура на чотирьох рівнях}: оператор~→~дрони~→~MQTT-брокер~→~хмарна аналітика. Усі компоненти задокументовані та мають посилання на академічні джерела — вимога, яку поставив автор на початку роботи і яку було повністю виконано.

  \item \textbf{Науково обґрунтований стек керування}: геометричний SE(3)-трекер~\cite{Lee2010, Mellinger2011} з квінтичними поліноміальними траєкторіями~\cite{Lynch2017}. Реалізований через підклас Pegasus~\texttt{NonlinearController}~\cite{Pegasus2024} без модифікації внутрішньої математики.

  \item \textbf{Покриття простору}: Boustrophedon Cellular Decomposition~\cite{Choset2000} — кожен дрон виконує один прямолінійний прохід вздовж осі~Y свого коридора. Поперечне відхилення $< 0{,}1$~м.

  \item \textbf{IoT-комунікація}: MQTT~v5.0~\cite{MQTT5} з розділенням топіків за критичністю (QoS~1 для команд і подій, QoS~0 для телеметрії). Телеметрія потрапляє у Grafana через Telegraf~→~InfluxDB~v2 без додаткової трансформації.

  \item \textbf{Естафетна передача місії}: при падінні заряду батареї нижче порогу система автоматично (без втручання оператора) передає другому дрону \emph{поточні світові координати} першого через той само MQTT-канал. Дрон-наступник приступає до сканування з відстані $< 0{,}5$~м від точки зупинки.

  \item \textbf{Grafana-дашборд}: 8 панелей реального часу — режим, позиція, швидкість, заряд батареї (\%), прогрес покриття (\%), інспекційні події.

  \item \textbf{Тестове покриття}: 49 одиничних тестів (pytest), всі проходять.

  \item \textbf{Переносимість}: система запускається як у фізичному симуляторі Isaac~Sim + Pegasus, так і у headless-режимі (\texttt{mock\_swarm}) без зміни коду місії.
\end{enumerate}

\subsection*{Напрями подальшого розвитку}

\begin{itemize}[itemsep=3pt]
  \item Заміна ground-truth позиції з PhysX на реальний VIO-стек (ORB-SLAM3~\cite{ORBSLAM3} або VINS-Fusion~\cite{VINSFusion}) для GPS-denied середовищ — архітектурно нульова зміна.
  \item Детектор аномалій на основі кадрів з наділових камер (YOLOv8 / SAM) для автоматичного виявлення дефектів стелажів.
  \item Автоматичний вибір дрона-наступника з пулу (не хардкодований \texttt{uav-02}) з оптимізацією за залишковим зарядом.
  \item Динамічне планування маршруту при виявленні перешкод (повне BCD~\cite{Choset2000} §3 або E-GTSP).
\end{itemize}
